\documentclass[12pt]{article}
\usepackage[utf8]{inputenc}
\usepackage[T1]{fontenc}
\usepackage{hyperref}
\title{Puroik Language Resources for K-12 Curriculum Development}
\author{Compiled by Claude Sonnet 5.5 (claude-sonnet-5-5) for Andr\'as Kornai}
\date{}
\begin{document}
\maketitle

\section{Introduction}
Puroik (also called Sulung, a name considered derogatory by neighbouring peoples) is spoken in Arunachal Pradesh, India, and in Lh\"unz\^e County, Tibet. Its classification is debated; Lieberherr places Bulu Puroik in the Kho-Bwa branch of Trans-Himalayan. Speaker estimates range from about 5,000--10,000 (Endangered Languages) to 20,000 (Wikipedia, 2011). Puroik is a dialect chain whose distant varieties are mutually unintelligible. Literacy is very low, reportedly about 2\%. The search found no school materials.

\section{Primary Educational Resources}
No primers, textbooks or school programmes in Puroik were found.

\section{Linguistic and Reference Materials}
\begin{itemize}
\item Ismael Lieberherr, \emph{A grammar of Bulu Puroik} (University of Bern, 2017), with 12 annotated texts and a Puroik--English dictionary in the appendix: \url{https://biblio.unibe.ch/download/eldiss/17lieberherr_i.pdf}
\item BORIS Theses record for the grammar: \url{https://boristheses.unibe.ch/2251}
\item Papers by Ismael Lieberherr: \url{https://unibe-ch.academia.edu/IsmaelLieberherr}
\item Puroik language (Wikipedia): \url{https://en.wikipedia.org/wiki/Puroik_language}
\item Endangered Languages Project entry: \url{https://endangeredlanguages.com/node/96719}. This entry cites Rai Soja, \emph{English-Puroik Dictionary} (Shillong, 2009) and Li Daqin, \emph{A study of Sulung} (Beijing, 2004). I did not find either online.
\end{itemize}

\section{Religious and Community Resources}
None found.

\section{Digital and Audio Resources}
None found beyond the annotated texts in the grammar.

\section{Curriculum Development Context}
The Bulu Puroik grammar and dictionary are the best base for materials, but they describe one variety within a dialect chain, and that variety was already receding at the time of fieldwork. Any curriculum should be developed with Puroik communities, and an agreed orthography and variety would need to be chosen first. I did not open the full PDF to check the dictionary's size.

\end{document}